\chapter{Scheduling and Backpressure}
\label{chap:scheduling}

The kernel schedules threads through a block--yield--wake cycle. Bounded ingress
propagates pressure, although not every internal queue has a hard memory bound.

\section{The kernel scheduler}

CharlotteOS uses a pluggable scheduler architecture. The system scheduler
(\texttt{SYSTEM\_SCHEDULER}) delegates to per-LP \texttt{LpScheduler} instances.
The reference implementation is \texttt{RoundRobin}:

\begin{itemize}
    \item A FIFO run queue per LP.
    \item A per-LP quantum timer (10 ms by default) that drives preemption.
    \item Operations: \texttt{spawn\_thread}, \texttt{block\_thread},
        \texttt{submit\_ready\_thread}, \texttt{yield\_lp}, \texttt{abort}.
\end{itemize}

The scheduler is pluggable: a different scheduling policy (EDF, CFS, fixed
priority) can be provided by implementing the \texttt{LpScheduler} trait. The
rest of the kernel interacts with the scheduler only through the abstract
interface.

\section{The block-yield-wake cycle}

This cycle underlies asynchronous operations. Every
async primitive in CharlotteOS --- \texttt{sleep()}, \texttt{wait()},
\texttt{CQ\_WAIT} --- is implemented through this cycle.

\subsection{Block}

A thread blocks by calling \texttt{block\_thread(tid, \&observable)}:

\begin{enumerate}
    \item The kernel marks the thread \texttt{Blocked} in the scheduler.
    \item The thread's \texttt{Waker} is registered as an observer on the
        observable event (a completion, a timer, an endpoint, or a CQ).
    \item The thread is removed from the LP's run queue.
    \item \texttt{yield\_lp()} saves the thread's register state and stack
        pointer to its \texttt{ThreadContext}.
    \item The scheduler selects the next ready thread (if any) and restores its
        context. The blocked thread is suspended at the point of the yield.
\end{enumerate}

For a timer sleep, steps 1--2 include creation of the wake source itself.
Publishing \texttt{Blocked} and inserting the timer into the local queue must
therefore be one local non-preemptible transaction. \texttt{sleep()} masks
local interrupts across both operations; otherwise a scheduler-quantum
interrupt can switch out the newly blocked thread before any event exists to
wake it. The previous interrupt state is restored before \texttt{yield\_lp()},
and no lock is held across the yield.

\subsection{Wake}

An event fires (timer expired, completion posted, message arrived):

\begin{enumerate}
    \item The event's observable fires, calling \texttt{Waker::notify(tid)}.
    \item \texttt{notify} calls \texttt{submit\_ready\_thread(tid)}, which marks
        the thread \texttt{Ready} and adds it to its affinity LP's run queue.
    \item If the target LP is idle, a reschedule IPI is sent to wake it.
        If the target LP is already running, no IPI is needed --- the thread
        will be picked up at the next scheduling point (quantum expiry or
        voluntary yield).
\end{enumerate}

\subsection{Resume}

At the next scheduling opportunity on the target LP:

\begin{enumerate}
    \item The scheduler selects the woken thread from the run queue.
    \item The thread's saved \texttt{ThreadContext} (registers, stack pointer)
        is restored.
    \item Execution resumes immediately after the \texttt{yield\_lp()} call that
        originally blocked the thread.
    \item From the thread's perspective, \texttt{block\_thread} returned normally
        and the event has fired.
\end{enumerate}

\section{Wake coalescing}

Wakes are \textbf{idempotent}: calling \texttt{submit\_ready\_thread} on an
already-\texttt{Ready} or \texttt{Running} thread is a no-op. This enables
coalescing:

\begin{itemize}
    \item If 10 messages arrive on an endpoint before the shard wakes, the kernel
        posts one wake, not 10.
    \item If 3 timer events fire before the LP processes the timer queue, the
        kernel delivers them in one batch during the next \texttt{process\_events}.
    \item The unified shard wait (\autoref{chap:sitas-executor}) aggregates multiple wake sources onto
        a single thread. Re-waking that thread is harmless.
\end{itemize}

The rule is: \textbf{enqueue before signaling.} The work is placed in a queue
(CQ ring, endpoint queue, timer queue) first; then a wake is delivered. The
woken thread finds the work when it drains the queue.

\section{The quantum timer}

Each LP's \texttt{RoundRobin} scheduler arms a periodic quantum timer (10 ms
default). When it fires:

\begin{enumerate}
    \item The timer IRQ handler sets a context-switch-pending flag.
    \item At the next \texttt{yield\_lp()} (or at the IRQ handler's tail), the
        scheduler checks the flag.
    \item If set and another thread is runnable, the current thread's context is
        saved and the next thread's is restored.
\end{enumerate}

An \texttt{armed} flag ensures exactly one quantum event is in flight at any time,
preventing unbounded timer-queue growth from manual yields.

\section{The deferred reaper}

A thread cannot free its own kernel stack (it is executing on it). When a thread
exits:

\begin{enumerate}
    \item \texttt{abort()} moves the dying thread from the thread table to a
        \texttt{DEAD\_THREADS} list. The thread is extracted from the table
        without being dropped.
    \item The scheduler performs a context switch to the next thread.
    \item \textbf{After} the switch (on the new thread's stack),
        \texttt{reap\_dead\_threads()} drops the dead threads, freeing their
        stacks and other resources.
\end{enumerate}

Deferred reclamation is required because threads own their stacks.

\section{Load rebalancing}

By default, threads are pinned to their affinity LP. Migration is opt-in:

\begin{itemize}
    \item A thread must be marked \texttt{migration\_safe}.
    \item Migration is rejected while the thread holds timer events, endpoint
        waiters, CQ registrations, or other LP-affine ownership constraints.
    \item Rebalancing uses a \textbf{sustained-imbalance window}: per-LP load is
        sampled over a runtime-adjustable interval rather than reacting to a
        single instantaneous queue-depth observation.
\end{itemize}

This conservative approach preserves cache warmth and eliminates migration races
(observer-not-found when a completion is posted on a different LP than the one
where the waiter was registered). Stable affinity is the correctness default;
migration is an opt-in optimization.

\subsection{Generation-safe retirement}

A numeric thread identifier is a recyclable table slot, not a permanent
identity. Every queued scheduler handle, waker, deferred verifier callback, and
asynchronous retirement request therefore carries the thread generation.
\texttt{abort\_thread\_generation} and the LP scheduler's removal operation
validate both values before changing state. A delayed cleanup request for an
exited thread consequently cannot abort a newer thread that reused the same
numeric identifier. The scheduler-lifecycle boot test exercises stale wake
rejection and remote owner-LP retirement.

\section{Synchronization and lock ordering}

Shared kernel state is synchronized by a small, deliberate set of primitives.
Their full semantics and the cross-subsystem ordering rules are maintained in
\repofile|docs/reference/locking.md|, which is authoritative over this
overview.

\subsection{Lock families}

\begin{itemize}
    \item \textbf{Interrupt-masking spin \texttt{Mutex}} guards the frame
        allocator, memory-object registry, address-space table, kernel address
        space, domain-authority table, address-space lifecycle, scratch-window
        cursor, and the global allocator. It attempts acquisition with
        maskable IRQs disabled and keeps them disabled while owned. A failed
        attempt restores the caller's interrupt state before spinning, then
        retries with IRQs masked. The ownership masking exists to prevent a
        preemption deadlock: these locks are taken from both preemptible kernel
        threads and synchronous EL0 exception paths, so a timer-preempted owner
        could otherwise be starved forever.

    \item \textbf{Interrupt-masking spin \texttt{RwLock}} is the workhorse. It
        protects the master thread table, the deferred-dead table, the system
        scheduler, the IPC registry, the completion registry, the userspace
        mailbox table, and each per-LP slot. Failed reader and writer attempts
        likewise restore \texttt{INT\_STATE} while contending. A
        \texttt{waiting\_writers} counter gives writers preference only during
        the masked atomic acquisition attempt; retaining that preference while
        waiting could deadlock a same-LP interrupt reader.

    \item \textbf{The external \texttt{spin} crate} is used only by the stack
        allocator, which wraps its arena mutex in explicit
        \texttt{mask\_interrupts!}/\texttt{unmask\_interrupts!} calls. Its
        guard-page map is touched only from thread context, never from IRQ
        context.

    \item \textbf{Lock-free structures} cover the rest: \texttt{ShardLocal} and
        \texttt{PerLp} for per-LP data, \texttt{ConcurrentQueue} for the
        IRQ-to-thread deferred-wake handoff, and plain atomics for the
        \texttt{on\_cpu} ownership handshake and generation counters.
\end{itemize}

A scheduler-blocking lock family exists under
\texttt{cpu/scheduler/sync} that parks a contended caller via
\texttt{block\_thread} instead of spinning. It currently has no callers and
does not participate in any ordering guarantee.

\subsection{Interrupt-masking discipline}

Both spin locks mask IRQs while owned, but leave IRQs in the caller's original
state while waiting for another LP. This is required for synchronous
TLB-shootdown rendezvous: a contender must remain able to acknowledge the
lock owner's IPI. Only the \texttt{RwLock} is nesting-aware. It
routes through a per-LP \texttt{INT\_STATE} record that counts save depth and
re-enables interrupts exactly once on the outermost restore. This lets a thread
take several different \texttt{RwLock}s, or a read after a read, without
prematurely re-enabling IRQs. The \texttt{Mutex} records its pre-acquire
interrupt bit on the lock object itself. Both guard types are non-sendable;
acquiring different mutexes still nests correctly, but re-acquiring the same
lock on one LP deadlocks.

\subsection{Lock ordering}

The scheduler's chain is fixed:

\begin{lstlisting}[style=rustCode]
SYSTEM_SCHEDULER.read() -> lp_scheduler.lock() -> MASTER_THREAD_TABLE.write()
\end{lstlisting}

Cross-subsystem rules hold today:

\begin{itemize}
    \item \textbf{Never take address-space or frame-allocator locks while the
        heap lock is held.} The heap's growth reserve is pre-mapped at boot so
        allocation can extend within mapped memory without those locks.
    \item \textbf{Never hold the memory-object registry across the
        address-space table lock.} Teardown takes table-then-frame-allocator
        while allocation takes allocator-then-registry; registry-then-table
        would close an AB-BC-CA cycle.
    \item \textbf{Dropping the memory-object registry for a bulk copy requires
        a shared-read copy pin.} While such a pin exists, writable CPU
        mappings, in-kernel writes, writable or exclusive DMA pins, write
        lends, moves, and ownership rollback are rejected. Owner teardown is
        deferred until both the DMA-pin and copy-pin counts reach zero. The
        final release removes the object under the registry lock and frees its
        frames only after dropping that lock.
    \item \textbf{Never block or yield while holding any lock.} A spin lock
        also masks IRQs, so parking the thread would abandon the lock and the
        LP could not schedule its successor. All guards are dropped before
        \texttt{switch\_ctx}.
    \item \textbf{IRQ context takes no locks and never blocks.} Interrupt
        delivery uses atomics and MMIO only; a deferred wake crosses into
        thread context through a lock-free queue.
\end{itemize}

\section{Backpressure and queue bounds}

Backpressure is a design principle. Endpoint and hardware rings are bounded,
but not every internal queue currently has a hard memory bound. In particular,
the non-lossy CQ overflow backlog may grow until its consumer drains it.

\begin{center}
\begin{tabular}{|l|l|l|}
\hline
\textbf{Layer} & \textbf{Bounded object} & \textbf{Backpressure signal} \\
\hline
Application request & Service endpoint queue & \texttt{QueueFull} \\
Service shard mailbox & \texttt{ShardSender<M>} & \texttt{Err(m)} returned \\
Driver endpoint/ring & NIC descriptor ring & Descriptor exhaustion \\
Kernel submission & Capability table & \texttt{WouldBlock} \\
Hardware & DMA descriptor ring & Ring full \\
\hline
\end{tabular}
\end{center}

Each layer has an explicit backpressure policy:

\begin{lstlisting}[style=rustCode]
enum SendPolicy {
    Try,               // Fail immediately if full.
    Wait,              // Wait for capacity.
    Deadline(Instant), // Wait until a deadline.
    // Merge a pending "data available" notification.
    DropIfCoalescible,
}
\end{lstlisting}

Rules:

\begin{itemize}
    \item Terminal completions are never dropped --- the CQ backlog
        guarantees this (\autoref{chap:completion-queues}).
    \item Coalescible notifications may be merged --- two ``data available''
        signals before the consumer wakes become one.
    \item Control messages normally fail or wait explicitly --- a service
        registration request that cannot be queued fails with \texttt{QueueFull}.
    \item Data-plane producers should stop before unbounded allocation ---
        a NIC driver should stop accepting TX packets when its descriptor ring
        is full.
    \item Emergency paths require reserved capacity --- the bounded IPI
        queue reserves slots for mandatory kernel messages (TLB shootdowns,
        scheduler commands), ensuring they cannot be starved by user work.
\end{itemize}

\section{Why scheduling and backpressure matter for critical software}

\begin{enumerate}
    \item The block-yield-wake cycle is exercised across subsystems --- the same
        mechanism handles sleep, IPC wait, CQ wait, and timer wait. A bug fixed
        in one path fixes all paths. The cycle is exercised by the async demo and
        by self-tests that call \texttt{wait()}; this is test evidence, not a
        formal proof.

    \item Idempotent wakes eliminate lost-wake races --- re-waking an
        already-ready thread is a no-op. This means the kernel can be aggressive
        about delivering wakes without worrying about double-delivery.

    \item Bounded ingress limits some exhaustion paths --- endpoint and
        descriptor-ring capacities reject excess work with backpressure.
        Resource quotas and a hard bound for all completion backlog memory are
        not yet implemented, so this is not a general denial-of-service
        guarantee.

    \item Stable LP affinity prevents migration races --- completion
        delivery is LP-local. A waker registered on LP 2 is signaled on LP 2.
        There is no window where a completion is posted on LP 1 while the thread
        is migrating to LP 2 and the observer is not found.

    \item The scheduling interface is pluggable --- \texttt{RoundRobin} is
        the implemented policy. Other policies would require their own
        implementation and validation behind the same interface.
\end{enumerate}

\section{Scheduler observability}

The prototype records integer running statistics for each thread at scheduler
context switches. A snapshot includes the thread identity and generation,
address-space owner, scheduler state, LP affinity, dispatch count, and the
count, minimum, maximum, total, and sum of squares of completed on-CPU slices.
The active slice start tick is reported separately. The counter frequency in
the snapshot header permits conversion from ticks to time.

The accumulator uses integer arithmetic and reports saturation. Mean and sample
variance are exported as exact rational components; standard deviation and
coefficient of variation are calculated in userspace. This avoids using
floating-point/SIMD state in the kernel. Accumulators are mergeable so future
hot-path metrics can be gathered per LP rather than through a global lock.

The \texttt{THREAD\_STATISTICS} syscall is address-space-scoped by default:
callers see only their own threads and receive the snapshot as a memory-object
capability. At boot the supervisor starts exactly one \texttt{observe} EL0
service and delegates a typed system-observer capability through its launch
vector. Presenting that capability authorizes a machine-wide snapshot; a
fabricated or wrong-type handle fails closed. The service publishes snapshots
through endpoint IPC and registers itself with the node name service.
Machine-wide observation is therefore explicit authority rather than an
ambient privilege.

An HTTP \texttt{GET /metrics} adapter is a reasonable userspace consumer of
this IPC protocol, but it is not implemented yet. The present TCP/IP service
does not supply the complete server-side bind/listen/accept path needed to host
it. See \repofile|docs/reference/observability.md| for the wire format, security boundary,
and suggested next metrics.

\endinput
